\subsubsection{The Heat Equation}\label{subsec:heat-equation}

The next elementary problem we examine is the \emph{heat equation}:
\begin{equation}\label{eq:heat-general}
u_t=\Delta u.
\end{equation}
Here $u$ is a real-valued function depending on "spatial" variables $\bm{x}\in\mathbb{R}^n$ and on "time" $t\in\mathbb{R}$, and the operator $\Delta$ is the Laplacian defined in \eqref{eq:laplace-general} which is assumed to act only on the spatial variables $(x_1,\ldots,x_n)$. (The reason for the quotation marks above is that in the next section we will describe the "type" of a differential equation in a way that is independent of any particular interpretation of the independent variables as spatial or temporal. However, even after we have done this, we will often lapse back to the terminology of space and time in order to draw analogies to the elementary Laplace, wave and heat equation described in this chapter.) As the name suggests, \eqref{eq:heat-general} describes the conduction of heat (with the dependent variable $u$ usually interpreted as temperature), but more generally it governs a range of physical phenomena described as \emph{diffusive}.

In discussing typical boundary conditions we confine ourselves to problems posed on a cylinder in space-time: $\Omega_t^+:=\{(\bm{x},t)\in\Omega\times(0,\infty)\}$ where $\Omega$ is a bounded domain in $\mathbb{R}^n$. Since the heat equation is first order in time we place one \emph{initial condition} on the solution. We let $\theta:\Omega\to\mathbb{R}$ be a given function and require
\begin{equation}\label{eq:heat-general-initial}
u(\bm{x},0)=\theta(\bm{x}).
\end{equation}
There are a variety of conditions typically posed on the boundary of the body.

\textbf{Temperature conditions.} Here we fix the dependent variable $u$ on some portion of the boundary.
\begin{equation}\label{eq:heat-temperature-boundary}
u(\bm{x},t)=f(\bm{x})
\end{equation}
for $\bm{x}\in\partial\Omega$ and $t\in(0,\infty)$. In problems of heat conduction, this corresponds to placing a portion of the boundary in contact with a constant temperature source (an ice bath, etc.). Of course, such conditions can be identified with Dirichlet conditions for Laplace's equation.

\textbf{Heat flux conditions.} Here we fix the normal derivative of $u$ on some portion of the boundary.
\begin{equation}\label{eq:heat-flux-boundary}
\dfrac{\partial u}{\partial n}(\bm{x},t)=g(\bm{x})
\end{equation}
for $\bm{x}\in\partial\Omega$ and $t\in(0,\infty)$, where $\bm{n}$ is the unit outward normal to $\partial\Omega$. A simplified version of Fourier's law of heat conduction says that the \emph{heat flux vector} $\bm{q}$ at a point $\bm{x}$ at time $t$ is given by
\begin{equation}\label{eq:fourier-heat-flux-law}
\bm{q}(\bm{x},t)=-\kappa\nabla u(\bm{x},t),
\end{equation}
where $\kappa$ is a positive constant called the \emph{thermal conductivity}. Thus, condition \eqref{eq:heat-flux-boundary} can be thought of as fixing the flow of heat through a portion of the boundary. (If $g=0$, we say that portion of the boundary is insulated.) The connection between heat flux conditions and Neumann conditions for Laplace's equation should be obvious.

\textbf{Linear radiation conditions.} Here we require
\begin{equation}\label{eq:linear-radiation-boundary}
-\dfrac{\partial u}{\partial n}(\bm{x},t)=\alpha u(\bm{x},t)
\end{equation}
for $\bm{x}\in\partial\Omega$ and $t\in(0,\infty)$, where $\alpha$ is a positive constant. This can be thought of as the linearization of \emph{Stefan's radiation law}
\begin{equation}\label{eq:stefan-radiation-law}
\bm{q}(\bm{x},t)\cdot\bm{n}(\bm{x})=\beta u^4(\bm{x},t)
\end{equation}
about a steady-state solution of the boundary-value problem. Stefan's law describes the loss of heat energy of a body through radiation into its surroundings.

\paragraph*{Solution by separation of variables}

As part of our review of elementary solution methods we now examine the solution of a one-dimensional heat conduction problem by the method of separation of variables. We consider the following \emph{initial/boundary-value problem}. Let
\[
D^+:=\{(x,t)\in\mathbb{R}^2\mid 0<x<1,\ 0<t<\infty\}.
\]
Find a function $u:\overline{D^+}\to\mathbb{R}$ that satisfies the differential equation
\begin{equation}\label{eq:heat-one-dimensional}
u_t=u_{xx}
\end{equation}
for $(x,t)\in D^+$, the initial condition
\begin{equation}\label{eq:heat-interval-initial}
u(x,0)=f(x)
\end{equation}
for $x\in(0,1)$, and the boundary conditions
\begin{equation}\label{eq:heat-left-boundary}
u(0,t)=0,
\end{equation}

\begin{equation}\label{eq:heat-right-boundary}
u(1,t)=0
\end{equation}
for $t>0$.

As before, we seek solutions of the form
\begin{equation}\label{eq:heat-separated-product}
u(x,t):=X(x)T(t).
\end{equation}
Plugging this into the differential equation \eqref{eq:heat-one-dimensional} gives us
\begin{equation}\label{eq:heat-product-substitution}
XT'=X''T.
\end{equation}
When $u$ is nonzero we get
\begin{equation}\label{eq:heat-separated-ratios}
\dfrac{T'(t)}{T(t)}=\dfrac{X''(x)}{X(x)}.
\end{equation}


Again, we make the argument that since the right side of the equation is independent of $t$ and the left side is independent of $x$, each side must be independent of both variables and hence must be equal to a constant. We then get the two ordinary differential equations
\begin{equation}\label{eq:heat-separated-time-equation}
T'=\lambda T,
\end{equation}

\begin{equation}\label{eq:heat-separated-space-equation}
X''=\lambda X.
\end{equation}


Solving these equations using \eqref{eq:heat-separated-product} gives us the following three-parameter family of solutions of \eqref{eq:heat-one-dimensional}:
\begin{equation}\label{eq:heat-three-parameter-family}
u(x,t)=e^{\lambda t}(Ae^{\sqrt{\lambda}x}+Be^{-\sqrt{\lambda}x}).
\end{equation}


Members of this family satisfy the boundary conditions \eqref{eq:heat-left-boundary} and \eqref{eq:heat-right-boundary} only when
\begin{equation}\label{eq:heat-boundary-eigenvalues}
\lambda=-n^2\pi^2,\qquad n=1,2,3,\ldots,
\end{equation}
in which case the collection of solutions reduces to
\begin{equation}\label{eq:heat-decaying-sine-modes}
u(x,t)=Ae^{-n^2\pi^2t}\sin n\pi x.
\end{equation}


To satisfy the initial condition, we use the Fourier expansion of the initial function $f$ via \eqref{eq:fourier-sine-series}, \eqref{eq:sine-coefficient-integral} to define
\begin{equation}\label{eq:heat-sine-series-solution}
u(x,t):=\sum_{n=1}^{\infty}A_ne^{-n^2\pi^2t}\sin n\pi x,
\end{equation}
where $A_n$ are equal to the Fourier coefficients of $f$. The questions that we asked at the end of Section~\ref{subsec:laplace-equation} concerning the series solution of Laplace's equation can be asked again about this solution of the heat equation. (Does the series converge? Does its limit satisfy the differential equation and the boundary conditions?) Once again, general answers await later chapters.

\begin{Example}
\label{ex:heat-tent-data}
We now formally compute a series solution for the initial boundary-value problem \eqref{eq:heat-one-dimensional}--\eqref{eq:heat-right-boundary}. In order to compare and contrast the solution we get here to the one obtained for Laplace's equation on page~\pageref{ex:tent-dirichlet-data} we use as our initial data the function $f$ defined in \eqref{eq:tent-boundary-data}. Since we have already computed the Fourier coefficients in \eqref{eq:tent-sine-coefficients} we can use these in \eqref{eq:heat-sine-series-solution} to define our solution.
\begin{equation}\label{eq:heat-tent-data-solution}
u(x,t):=\sum_{k=0}^{\infty}\dfrac{4(-1)^k}{(2k+1)^2\pi^2}e^{-(2k+1)^2\pi^2t}\sin(2k+1)\pi x.
\end{equation}

\end{Example}

\paragraph*{Instability of backwards heat equation}

In this section we consider the following problem. Let
\[
D^-:=\{(x,t)\in\mathbb{R}^2\mid 0<x<1,\ -\infty<t<0\}.
\]
We wish to find a function $u:\overline{D^-}\to\mathbb{R}$ that satisfies the heat equation \eqref{eq:heat-one-dimensional} in $D^-$, the boundary conditions \eqref{eq:heat-left-boundary} and \eqref{eq:heat-right-boundary} for $t<0$, and the \emph{terminal condition}
\begin{equation}\label{eq:heat-terminal-condition}
u(x,0)=f(x)
\end{equation}
for $x\in(0,1)$.

The problem is often transformed using the change of variables
\begin{equation}\label{eq:heat-time-reversal}
\bar t=-t.
\end{equation}
Under this transformation we seek to solve the differential equation
\begin{equation}\label{eq:backwards-heat-equation}
u_{\bar t}=-u_{xx}
\end{equation}
for $(x,\bar t)\in D^+$. The boundary conditions \eqref{eq:heat-left-boundary} and \eqref{eq:heat-right-boundary} are now applied for positive $t$, and the terminal condition \eqref{eq:heat-terminal-condition} remains unchanged but is now thought of as an initial condition. This version of the problem is known as the \emph{backwards heat equation}. The name stems, of course, from the fact that in the original formulation we are trying to solve the heat equation for times \emph{before} the terminal condition \eqref{eq:heat-terminal-condition}.

The formal separation of variables solution for the backwards heat equation is
\begin{equation}\label{eq:backwards-heat-series}
u(x,t):=\sum_{n=1}^{\infty}A_ne^{n^2\pi^2t}\sin n\pi x,
\end{equation}
where $A_n$ are the Fourier coefficients of $f$. The derivation is almost exactly the same as before with the sole difference that we end up with exponentials that \emph{grow rather than decay with time}. As a result we can get the following instability result that states that we can find initial data for which the solution of the backwards heat equation blows up as quickly as desired.

\begin{Lemma}
\label{lem:backwards-heat-instability}
For any $T>0$, $M>0$, and $\epsilon>0$ there exists an initial function $f$ satisfying
\begin{equation}\label{eq:backwards-small-initial-norm}
\|f\|_{C([0,1])}=\epsilon
\end{equation}
such that the backwards heat conduction problem has a separation of variables solution $u(x,t)$ defined by \eqref{eq:backwards-heat-series} that satisfies
\begin{equation}\label{eq:backwards-large-terminal-norm}
\|u(\cdot,T)\|_{C([0,1])}\geq M.
\end{equation}

\end{Lemma}

Here $\|f\|_{C([0,1])}$ denotes the max of $|f|$ on $[0,1]$.

\begin{Proof}
Choose $n$ sufficiently large so that
\begin{equation}\label{eq:backwards-frequency-choice}
n^2\pi^2T\geq\ln\left(\dfrac M\epsilon\right).
\end{equation}
Then if
\begin{equation}\label{eq:backwards-sine-initial-data}
f(x)=\epsilon\sin n\pi x,
\end{equation}
the solution
\begin{equation}\label{eq:backwards-growing-sine-solution}
u(x,t)=\epsilon e^{n^2\pi^2t}\sin n\pi x
\end{equation}
satisfies the requirement.
\end{Proof}

\paragraph*{Energy inequality}

We end our discussion of the heat equation by proving an \emph{energy estimate}. The simple estimate we derive in this section should act as a prototype for estimates that we will derive in later chapters. We will show the following.

\begin{Lemma}
\label{lem:heat-energy-decay}
Let $u:\overline{D^+}\to\mathbb{R}$ be a $C^2$ solution of the heat equation \eqref{eq:heat-one-dimensional} satisfying the boundary conditions \eqref{eq:heat-left-boundary} and \eqref{eq:heat-right-boundary}. Then for any $t_1\geq t_0\geq0$, the solution $u$ satisfies
\begin{equation}\label{eq:heat-energy-decay}
\int_0^1u^2(x,t_1)\,\mathrm{d}x\leq\int_0^1u^2(x,t_0)\,\mathrm{d}x.
\end{equation}

\end{Lemma}

In the language of Chapter~\ref{chap:function-spaces}, for any solution of the heat equation satisfying the given boundary conditions, the $L^2$ norm (in space) decreases with time.

\begin{Proof}
We first use the heat equation to derive the following differential identity for $u$.
\begin{equation}\label{eq:heat-pointwise-energy-identity}
\begin{aligned}
\dfrac{\mathrm{d}}{\mathrm{d}t}u^2&=2uu_t\\
&=2uu_{xx}\\
&=2(uu_x)_x-2u_x^2.
\end{aligned}
\end{equation}
Integrating both sides of this identity with respect to $x$ gives us
\begin{equation}\label{eq:heat-integrated-energy-identity}
\begin{aligned}
\dfrac12\int_0^1(u^2(x,t))_t\,\mathrm{d}x
&=u(1,t)u_x(1,t)-u(0,t)u_x(0,t)\\
&\quad-\int_0^1u_x^2(x,t)\,\mathrm{d}x.
\end{aligned}
\end{equation}
We now use the boundary conditions to eliminate the boundary terms in the equation above and integrate the result with respect to time. After changing the order of integration on the left side we get
\begin{equation}\label{eq:heat-energy-dissipation}
\begin{aligned}
\int_0^1u^2(x,t_1)\,\mathrm{d}x-\int_0^1u^2(x,t_0)\,\mathrm{d}x
&=\int_0^1\int_{t_0}^{t_1}(u^2)_t\,\mathrm{d}t\,\mathrm{d}x\\
&=-2\int_0^1\int_{t_0}^{t_1}u_x^2\,\mathrm{d}t\,\mathrm{d}x\leq0.
\end{aligned}
\end{equation}
This completes the proof.
\end{Proof}

\paragraph*{Problems}

\begin{enumerate}[label=\thesection.\arabic*.,ref=\thesection.\arabic*,start=21,series=book-problems]

\item\label{ex:heat-three-boundary-problems} Solve the one-dimensional heat equation via separation of variables for the following boundary conditions:\begin{enumerate}[label=(\alph*)]
\item $\begin{aligned}u(x,0)&=x^2-x,\\u_x(0,t)&=0,\\u_x(1,t)&=0.\end{aligned}$
\item $\begin{aligned}u(x,0)&=\sin\pi x,\\u(0,t)&=0,\\u_x(1,t)&=0.\end{aligned}$
\item $\begin{aligned}
u(x,0)&=\begin{cases}x,&x\in[0,1/2)\\1-x,&x\in[1/2,1],\end{cases}\\
u(0,t)&=0,\\u(1,t)&=0.
\end{aligned}$
\end{enumerate}

\item\label{ex:heat-physical-time-scaling} In a typical physical problem in heat conduction, one studies the differential equation
\[
c\rho u_t=\kappa\Delta u
\]
where $c$ is the \emph{specific heat}, $\rho$ is the \emph{density}, and $\kappa$ is the thermal conductivity of the medium under consideration. If $c$, $\rho$, and $\kappa$ are constants, show that there is a linear change in time scale $\bar t=\gamma t$ that transforms the differential equation above into \eqref{eq:heat-general}.

\item\label{ex:duhamel-heat-pulse} Suppose $f:\overline{D^+}\to\mathbb{R}$ is continuous and $u:\overline{D^+}\to\mathbb{R}$ is a solution of the following nonhomogeneous initial/boundary-value problem:
\[
\begin{gathered}
u_t(x,t)-u_{xx}(x,t)=f(x,t),\quad(x,t)\in D^+,\\
u(x,0)=0,\quad x\in[0,1],\\
u(0,t)=u(1,t)=0,\quad t\in[0,\infty).
\end{gathered}
\]
Now, for each $\tau\in[0,\infty)$, let $w(x,t,\tau)$ be the solution of the following \emph{pulse problem}:
\[
\begin{gathered}
w_t-w_{xx}=0,\quad(x,t)\in(0,1)\times(\tau,\infty),\\
w(x,\tau,\tau)=f(x,\tau),\quad x\in[0,1],\\
w(0,t,\tau)=w(1,t,\tau)=0,\quad t\in[\tau,\infty).
\end{gathered}
\]
Show that $u$ and $w$ satisfy the relation
\[
u(x,t)=\int_0^tw(x,t,\tau)\,\mathrm{d}\tau.
\]
This and similar methods of relating nonhomogeneous PDEs with homogeneous initial conditions to homogeneous PDEs with nonhomogeneous initial conditions are known as Duhamel's principle.

\item\label{ex:heat-step-cauchy-problem} Solve the Cauchy problem
\[
u_t=u_{xx},\qquad u(x,0)=\begin{cases}0,&x<0\\1,&x>0.\end{cases}
\]
Hint: Seek a solution in the form $u(x,t)=\phi(x/\sqrt t)$.

\end{enumerate}
